
\section{Цель работы}

Целью данной лабораторной работы является освоение методов принятия решений на основе нечётких отношений предпочтения. В отличие от классических задач, рассмотренных в предыдущих работах, здесь предпочтения лица, принимающего решения (ЛПР), задаются не точными числовыми оценками, а степенями уверенности в превосходстве одной альтернативы над другой. Такие степени принимают значения из отрезка $[0,1]$, что позволяет естественным образом описывать неполную, неточную и противоречивую экспертную информацию.

В рамках работы требуется освоить построение пересечения нечётких отношений предпочтения, научиться вычислять транспонированное отношение и строгое отношение предпочтения, а также определять степени недоминируемости альтернатив. Кроме того, необходимо освоить операцию взвешенной свёртки нескольких отношений предпочтения с учётом весовых коэффициентов $\lambda_1, \ldots, \lambda_7$ и научиться объединять результаты, полученные для пересечения и для взвешенной свёртки, в итоговую оценку альтернатив.

Практической целью работы является решение конкретной задачи для варианта 4: по заданным матрицам $R_1, \ldots, R_7$ и весовым коэффициентам найти альтернативу, обладающую наибольшей степенью недоминируемости, и содержательно проинтерпретировать полученный результат.

\clearpage

\section{Теоретическая часть}

Пусть $X = \{x_1, \ldots, x_n\}$ --- конечное множество альтернатив. Нечётким отношением предпочтения на этом множестве называется нечёткое множество, заданное на декартовом произведении $X \times X$. Иными словами, это матрица $R = (r_{ij})$ размерности $n \times n$, каждый элемент которой характеризует степень уверенности в том, что $x_i$ предпочтительнее $x_j$. Значения $r_{ij}$ лежат в отрезке $[0,1]$ и отражают силу превосходства: $0$ означает его отсутствие, $0.5$ --- неопределённость предпочтения, $1$ --- максимальное превосходство.

В общем случае нечёткое отношение предпочтения не обязано быть ни симметричным, ни обратно симметричным: величины $r_{ij}$ и $r_{ji}$ задаются экспертами независимо и могут не дополнять друг друга до единицы. Поэтому для корректного анализа строится ряд производных отношений.

\textbf{Пересечение нечётких отношений.} Если имеется несколько нечётких отношений предпочтения $R_1, R_2, \ldots, R_m$, каждое из которых отражает мнение отдельного эксперта или отдельного критерия, то их совместное мнение описывается пересечением
\[
Q_1 = R_1 \cap R_2 \cap \ldots \cap R_m.
\]
Для нечётких множеств операция пересечения определяется покоординатным минимумом:
\[
q_{ij} = \min_{k=1,\ldots,m} r_{ij}^{(k)}, \quad i,j = 1,\ldots,n.
\]
Полученное отношение $Q_1$ сохраняет только те предпочтения, которые одновременно присутствуют во всех исходных отношениях. Чем больше экспертов согласны с тем, что $x_i$ превосходит $x_j$, тем выше соответствующее значение $q_{ij}$. Таким образом, пересечение является консервативной операцией: оно не допускает, чтобы единственный эксперт «навязал» своё мнение всей группе.

\textbf{Транспонированное отношение.} Для отношения $Q = (q_{ij})$ транспонированным называется отношение $Q^{*} = (q_{ij}^{*})$, элементы которого определяются как
\[
q_{ij}^{*} = q_{ji}.
\]
Транспонированное отношение описывает те же предпочтения, но с противоположной точки зрения: элемент $q_{ij}^{*}$ показывает, в какой степени $x_i$ уступает $x_j$. Иными словами, $Q^{*}$ получается из $Q$ зеркальным отражением относительно главной диагонали.

\textbf{Строгое отношение предпочтения.} Разность между отношением $Q$ и его транспонированным отношением позволяет выделить «чистое» превосходство, свободное от взаимных уступок. Строгое отношение предпочтения (н.о.с.п.) $Q_s$ определяется как
\[
Q_s = Q \setminus Q^{*},
\]
что для нечётких отношений раскрывается формулой
\[
q^{s}_{ij} = \max\bigl(0,\; q_{ij} - q_{ji}\bigr), \quad i,j = 1,\ldots,n.
\]
Если $q_{ij} > q_{ji}$, то альтернатива $x_i$ превосходит $x_j$ в степени $q_{ij} - q_{ji}$; если же $q_{ij} \le q_{ji}$, то превосходство отсутствует и соответствующая степень равна нулю. Диагональные элементы строгого отношения всегда равны нулю, поскольку $q_{ii} - q_{ii} = 0$.

\textbf{Степени недоминируемости.} Для выбора наилучшей альтернативы используется понятие степени недоминируемости. Сначала для каждой альтернативы $x_i$ определяется максимальная степень, с которой какая-либо другая альтернатива превосходит её:
\[
d_i = \max_{j=1,\ldots,n} q^{s}_{ji}.
\]
Здесь индекс $ji$ означает, что рассматривается столбец $i$ строгого отношения, то есть насколько другие альтернативы «подавляют» $x_i$. Степень недоминируемости определяется как дополнение этой величины до единицы:
\[
\mu_{\text{н.д.},Q}(x_i) = 1 - d_i = 1 - \max_{j=1,\ldots,n} q^{s}_{ji}.
\]
Альтернатива, которую никто не превосходит, получает степень недоминируемости, равную единице; альтернатива, которая сильно подавляется другими, получает значение, близкое к нулю. Чем выше $\mu_{\text{н.д.},Q}(x_i)$, тем более предпочтительной является альтернатива.

\textbf{Взвешенная свёртка отношений.} Пересечение учитывает только согласованные предпочтения и не различает экспертов по их компетентности. Если эксперты или критерии имеют различную важность, задаваемую весами $\lambda_1, \ldots, \lambda_m$, то строится взвешенная свёртка
\[
Q_2 = \lambda_1 R_1 + \lambda_2 R_2 + \ldots + \lambda_m R_m,
\]
элементы которой вычисляются по формуле
\[
q^{(2)}_{ij} = \sum_{k=1}^{m} \lambda_k r_{ij}^{(k)}, \quad i,j = 1,\ldots,n.
\]
Весовые коэффициенты предполагаются нормированными:
\[
\sum_{k=1}^{m} \lambda_k = 1, \qquad \lambda_k \ge 0.
\]
В отличие от пересечения, взвешенная свёртка является компромиссной операцией: она учитывает мнение всех экспертов, но вклад каждого пропорционален его весу. В результате элементы $Q_2$ могут быть больше, чем соответствующие элементы $Q_1$, поскольку высокие оценки одних экспертов частично компенсируют низкие оценки других.

К отношению $Q_2$ применяется та же процедура, что и к $Q_1$: вычисляются транспонированное отношение $Q_2^{*}$, строгое отношение $Q_{2s} = Q_2 \setminus Q_2^{*}$ и степени недоминируемости $\mu_{\text{н.д.},Q_2}(x_i)$.

\textbf{Итоговое правило выбора.} Два подхода — пересечение и взвешенная свёртка — могут давать различные рекомендации. Чтобы получить устойчивое решение, итоговая степень недоминируемости альтернативы определяется как минимум из двух полученных значений:
\[
\mu_{\text{н.д.}}(x_i) = \min\Bigl(\mu_{\text{н.д.},Q_1}(x_i),\;\mu_{\text{н.д.},Q_2}(x_i)\Bigr).
\]
Оптимальными считаются альтернативы, для которых эта величина максимальна:
\[
x^{*} = \arg\max_{x_i \in X} \mu_{\text{н.д.}}(x_i).
\]
Использование минимума означает, что альтернатива должна быть хороша одновременно и по консервативному, и по компромиссному критерию; это повышает надёжность итогового выбора.

% \clearpage

\section{Практическая часть}

В соответствии с вариантом 4 на множестве из шести альтернатив $X = \{x_1, \ldots, x_6\}$ заданы семь нечётких отношений предпочтения $R_1, \ldots, R_7$ и следующие весовые коэффициенты:
\[
\lambda_1 = 0.1, \quad \lambda_2 = 0.2, \quad \lambda_3 = 0.1, \quad \lambda_4 = 0.1, \quad \lambda_5 = 0.3, \quad \lambda_6 = 0.1, \quad \lambda_7 = 0.1.
\]
Сумма весов равна единице:
\[
0.1 + 0.2 + 0.1 + 0.1 + 0.3 + 0.1 + 0.1 = 1.
\]
Исходные матрицы отношений предпочтения имеют вид:

\begin{multicols}{2}
\[
R_1 = \begin{pmatrix}
1 & 0.3 & 0.8 & 0.5 & 0.5 & 0.6 \\
0.4 & 1 & 0.4 & 0.9 & 0.2 & 0.2 \\
0.7 & 0.7 & 1 & 0.8 & 0.7 & 0.7 \\
0.1 & 0.3 & 0.7 & 1 & 0.4 & 0.7 \\
0.8 & 0.9 & 0.7 & 0.4 & 1 & 0.4 \\
0.8 & 0.4 & 0.2 & 0.3 & 0.2 & 1
\end{pmatrix}
\]

\[
R_2 = \begin{pmatrix}
1 & 0.8 & 0.1 & 0.6 & 0.9 & 0.3 \\
0.4 & 1 & 0.5 & 0.3 & 0.7 & 0.4 \\
0.2 & 0.7 & 1 & 0.7 & 0.7 & 0.6 \\
0.7 & 0.8 & 0.8 & 1 & 0.3 & 0.3 \\
0.3 & 0.3 & 0.6 & 0.6 & 1 & 0.2 \\
0.3 & 0.6 & 0.8 & 0.6 & 0.1 & 1
\end{pmatrix}
\]

\[
R_3 = \begin{pmatrix}
1 & 0.4 & 0.7 & 0.6 & 0.3 & 0.3 \\
0.3 & 1 & 0.3 & 0.1 & 0.5 & 0.7 \\
0.3 & 0.8 & 1 & 0.3 & 0.9 & 0.2 \\
0.2 & 0.7 & 0.3 & 1 & 0.7 & 0.8 \\
0.3 & 0.8 & 0.8 & 0.6 & 1 & 0.4 \\
0.4 & 0.7 & 0.2 & 0.9 & 0.4 & 1
\end{pmatrix}
\]

\[
R_4 = \begin{pmatrix}
1 & 0.6 & 0.3 & 0.7 & 0.7 & 0.2 \\
0.8 & 1 & 0.7 & 0.5 & 0.3 & 0.3 \\
0.8 & 0.1 & 1 & 0.7 & 0.1 & 0.7 \\
0.3 & 0.8 & 0.6 & 1 & 0.2 & 0.2 \\
0.5 & 0.5 & 0.8 & 0.7 & 1 & 0.2 \\
0.8 & 0.5 & 0.2 & 0.6 & 0.1 & 1
\end{pmatrix}
\]

\[
R_5 = \begin{pmatrix}
1 & 0.2 & 0.5 & 0.5 & 0.5 & 0.2 \\
0.6 & 1 & 0.6 & 0.1 & 0.8 & 0.9 \\
0.7 & 0.8 & 1 & 0.3 & 0.8 & 0.4 \\
0.4 & 0.6 & 0.6 & 1 & 0.4 & 0.8 \\
0.2 & 0.9 & 0.1 & 0.2 & 1 & 0.7 \\
0.4 & 0.8 & 0.3 & 0.8 & 0.5 & 1
\end{pmatrix}
\]

\[
R_6 = \begin{pmatrix}
1 & 0.7 & 0.7 & 0.4 & 0.6 & 0.4 \\
0.2 & 1 & 0.2 & 0.4 & 0.1 & 0.4 \\
0.3 & 0.3 & 1 & 0.8 & 0.2 & 0.5 \\
0.8 & 0.5 & 0.7 & 1 & 0.8 & 0.3 \\
0.7 & 0.3 & 0.6 & 0.6 & 1 & 0.9 \\
0.4 & 0.7 & 0.6 & 0.7 & 0.8 & 1
\end{pmatrix}
\]

\[
R_7 = \begin{pmatrix}
1 & 0.7 & 0.8 & 0.8 & 0.3 & 0.7 \\
0.2 & 1 & 0.7 & 0.5 & 0.6 & 0.8 \\
0.5 & 0.8 & 1 & 0.8 & 0.7 & 0.9 \\
0.3 & 0.9 & 0.6 & 1 & 0.4 & 0.9 \\
0.4 & 0.5 & 0.4 & 0.7 & 1 & 0.1 \\
0.7 & 0.3 & 0.4 & 0.8 & 0.6 & 1
\end{pmatrix}
\]
\end{multicols}

% \clearpage

\textbf{Этап 1. Построение пересечения $Q_1$.}

Пересечение строится покоординатным минимумом по всем семи матрицам:
\[
q_{ij} = \min\bigl(r^{(1)}_{ij}, r^{(2)}_{ij}, \ldots, r^{(7)}_{ij}\bigr).
\]
Покажем вычисление двух элементов. Для пары $(x_1, x_2)$ имеем:
\[
q_{12} = \min(0.3,\; 0.8,\; 0.4,\; 0.6,\; 0.2,\; 0.7,\; 0.7) = 0.2.
\]
Для пары $(x_3, x_4)$:
\[
q_{34} = \min(0.8,\; 0.7,\; 0.3,\; 0.7,\; 0.3,\; 0.8,\; 0.8) = 0.3.
\]
Аналогично вычисляются все остальные элементы. В результате получаем отношение
\[
Q_1 = \begin{pmatrix}
1 & 0.2 & 0.1 & 0.4 & 0.3 & 0.2 \\
0.2 & 1 & 0.2 & 0.1 & 0.1 & 0.2 \\
0.2 & 0.1 & 1 & 0.3 & 0.1 & 0.2 \\
0.1 & 0.3 & 0.3 & 1 & 0.2 & 0.2 \\
0.2 & 0.3 & 0.1 & 0.2 & 1 & 0.1 \\
0.3 & 0.3 & 0.2 & 0.3 & 0.1 & 1
\end{pmatrix}
\]

% \clearpage

\textbf{Этап 2. Транспонированное отношение $Q_1^{*}$ и строгое отношение $Q_{1s}$.}

Транспонированное отношение получается зеркальным отражением матрицы $Q_1$ относительно главной диагонали:
\[
Q_1^{*} = \begin{pmatrix}
1 & 0.2 & 0.2 & 0.1 & 0.2 & 0.3 \\
0.2 & 1 & 0.1 & 0.3 & 0.3 & 0.3 \\
0.1 & 0.2 & 1 & 0.3 & 0.1 & 0.2 \\
0.4 & 0.1 & 0.3 & 1 & 0.2 & 0.3 \\
0.3 & 0.1 & 0.1 & 0.2 & 1 & 0.1 \\
0.2 & 0.2 & 0.2 & 0.2 & 0.1 & 1
\end{pmatrix}
\]
Строгое отношение предпочтения вычисляется по формуле
\[
q^{s}_{ij} = \max\bigl(0,\; q_{ij} - q_{ji}\bigr).
\]
Например, для пары $(x_1, x_4)$:
\[
q^{s}_{14} = \max(0,\; 0.4 - 0.1) = 0.3,
\]
а для пары $(x_2, x_3)$:
\[
q^{s}_{23} = \max(0,\; 0.2 - 0.1) = 0.1.
\]
Если $q_{ij} \le q_{ji}$, соответствующий элемент обнуляется:
\[
q^{s}_{21} = \max(0,\,0.2-0.2) = 0.
\]
Полностью строгое отношение имеет вид
\[
Q_{1s} = \begin{pmatrix}
0 & 0 & 0 & 0.3 & 0.1 & 0 \\
0 & 0 & 0.1 & 0 & 0 & 0 \\
0.1 & 0 & 0 & 0 & 0 & 0 \\
0 & 0.2 & 0 & 0 & 0 & 0 \\
0 & 0.2 & 0 & 0 & 0 & 0 \\
0.1 & 0.1 & 0 & 0.1 & 0 & 0
\end{pmatrix}
\]

% \clearpage

\textbf{Этап 3. Степени недоминируемости по отношению $Q_1$.}

Для каждой альтернативы находим максимальное значение в соответствующем столбце матрицы $Q_{1s}$ и вычитаем его из единицы:
\[
\mu_{\text{н.д.},Q_1}(x_i) = 1 - \max_{j} q^{s}_{ji}.
\]
Максимумы по столбцам равны соответственно
\[
(0.1,\; 0.2,\; 0.1,\; 0.3,\; 0.1,\; 0),
\]
поэтому
\[
\mu_{\text{н.д.},Q_1} = (0.9,\; 0.8,\; 0.9,\; 0.7,\; 0.9,\; 1.0).
\]
Таким образом, по консервативному критерию (пересечению) наибольшую степень недоминируемости имеет альтернатива $x_6$: её не превосходит ни одна другая альтернатива, и $\mu_{\text{н.д.},Q_1}(x_6) = 1$. Альтернативы $x_1$, $x_3$ и $x_5$ делят второе место со значением $0.9$, а наихудшей является $x_4$ со значением $0.7$.

% \clearpage

\textbf{Этап 4. Взвешенная свёртка $Q_2$.}

Взвешенная свёртка вычисляется как линейная комбинация исходных отношений:
\[
q^{(2)}_{ij} = \sum_{k=1}^{7} \lambda_k r^{(k)}_{ij}.
\]
Покажем подробно вычисление элемента $q^{(2)}_{12}$:
\[
\begin{aligned}
q^{(2)}_{12} &= 0.1 \cdot 0.3 + 0.2 \cdot 0.8 + 0.1 \cdot 0.4 + 0.1 \cdot 0.6 + \\
&\quad + 0.3 \cdot 0.2 + 0.1 \cdot 0.7 + 0.1 \cdot 0.7 = \\
&= 0.03 + 0.16 + 0.04 + 0.06 + 0.06 + 0.07 + 0.07 = 0.49.
\end{aligned}
\]
Аналогично для $q^{(2)}_{13}$:
\[
\begin{aligned}
q^{(2)}_{13} &= 0.1 \cdot 0.8 + 0.2 \cdot 0.1 + 0.1 \cdot 0.7 + 0.1 \cdot 0.3 +\\
&\quad + 0.3 \cdot 0.5 + 0.1 \cdot 0.7 + 0.1 \cdot 0.8 = \\
&= 0.08 + 0.02 + 0.07 + 0.03 + 0.15 + 0.07 + 0.08 = 0.50.
\end{aligned}
\]
Выполнив такие вычисления для всех пар, получаем
\[
Q_2 = \begin{pmatrix}
1 & 0.49 & 0.50 & 0.57 & 0.57 & 0.34 \\
0.45 & 1 & 0.51 & 0.33 & 0.55 & 0.59 \\
0.51 & 0.65 & 1 & 0.57 & 0.64 & 0.54 \\
0.43 & 0.66 & 0.63 & 1 & 0.43 & 0.59 \\
0.39 & 0.63 & 0.48 & 0.48 & 1 & 0.45 \\
0.49 & 0.62 & 0.41 & 0.69 & 0.38 & 1
\end{pmatrix}
\]

% \clearpage

\textbf{Этап 5. Транспонированное и строгое отношения для $Q_2$.}

Транспонированное отношение имеет вид
\[
Q_2^{*} = \begin{pmatrix}
1 & 0.45 & 0.51 & 0.43 & 0.39 & 0.49 \\
0.49 & 1 & 0.65 & 0.66 & 0.63 & 0.62 \\
0.50 & 0.51 & 1 & 0.63 & 0.48 & 0.41 \\
0.57 & 0.33 & 0.57 & 1 & 0.48 & 0.69 \\
0.57 & 0.55 & 0.64 & 0.43 & 1 & 0.38 \\
0.34 & 0.59 & 0.54 & 0.59 & 0.45 & 1
\end{pmatrix}
\]
Строгое отношение $Q_{2s} = Q_2 \setminus Q_2^{*}$ вычисляется по той же формуле $q^{s}_{ij} = \max(0,\,q_{ij}-q_{ji})$. Например:
\[
q^{s}_{14} = \max(0,\; 0.57 - 0.43) = 0.14,
\]
\[
q^{s}_{42} = \max(0,\; 0.66 - 0.33) = 0.33,
\]
\[
q^{s}_{15} = \max(0,\; 0.57 - 0.39) = 0.18.
\]
В случае $q_{ij} < q_{ji}$ элемент равен нулю, например $q^{s}_{21} = \max(0,\,0.45-0.49) = 0$. В результате
\[
Q_{2s} = \begin{pmatrix}
0 & 0.04 & 0 & 0.14 & 0.18 & 0 \\
0 & 0 & 0 & 0 & 0 & 0 \\
0.01 & 0.14 & 0 & 0 & 0.16 & 0.13 \\
0 & 0.33 & 0.06 & 0 & 0 & 0 \\
0 & 0.08 & 0 & 0.05 & 0 & 0.07 \\
0.15 & 0.03 & 0 & 0.10 & 0 & 0
\end{pmatrix}
\]

% \clearpage

\textbf{Этап 6. Степени недоминируемости по отношению $Q_2$.}

Максимумы по столбцам матрицы $Q_{2s}$ равны
\[
(0.15,\; 0.33,\; 0.06,\; 0.14,\; 0.18,\; 0.13),
\]
поэтому
\[
\mu_{\text{н.д.},Q_2} = (0.85,\; 0.67,\; 0.94,\; 0.86,\; 0.82,\; 0.87).
\]
В отличие от пересечения, взвешенная свёртка отдаёт первое место альтернативе $x_3$ со значением $0.94$. Это объясняется тем, что при компромиссном усреднении высокие оценки, выставленные экспертами с большими весами, частично компенсируют низкие оценки остальных. Альтернатива $x_2$ получает наименьшее значение $0.67$, поскольку в столбце $x_2$ строгого отношения присутствует большой элемент $q^{s}_{42}=0.33$, отражающий сильное превосходство $x_4$ над $x_2$.

% \clearpage

\textbf{Этап 7. Итоговые степени недоминируемости и выбор оптимальной альтернативы.}

Итоговая степень недоминируемости вычисляется как минимум двух полученных значений:
\[
\mu_{\text{н.д.}}(x_i) = \min\bigl(\mu_{\text{н.д.},Q_1}(x_i),\; \mu_{\text{н.д.},Q_2}(x_i)\bigr).
\]
Результаты сведены в таблицу:

\begin{center}
\begin{tabular}{|c|c|c|c|}
\hline
Альтернатива & $\mu_{\text{н.д.},Q_1}$ & $\mu_{\text{н.д.},Q_2}$ & $\mu_{\text{н.д.}}$ \\
\hline
$x_1$ & 0.9 & 0.85 & 0.85 \\
$x_2$ & 0.8 & 0.67 & 0.67 \\
$x_3$ & 0.9 & 0.94 & \textbf{0.9} \\
$x_4$ & 0.7 & 0.86 & 0.7 \\
$x_5$ & 0.9 & 0.82 & 0.82 \\
$x_6$ & 1.0 & 0.87 & 0.87 \\
\hline
\end{tabular}
\end{center}

Максимальное значение итоговой степени недоминируемости равно
\[
\max_i \mu_{\text{н.д.}}(x_i) = 0.9,
\]
и достигается оно у альтернативы $x_3$. Следовательно, оптимальной альтернативой является $x_3$, для которой
\[
\mu_{\text{н.д.}}(x_3) = 0.9.
\]

Отметим, что альтернатива $x_6$ имела максимальную степень недоминируемости по пересечению ($\mu = 1$), однако по взвешенной свёртке она уступает $x_3$ ($0.87$ против $0.94$). Именно операция взятия минимума не позволяет выбрать $x_6$: её преимущество оказывается неустойчивым при переходе к компромиссному критерию. Альтернатива $x_3$, напротив, показывает высокие результаты в обоих случаях ($0.9$ и $0.94$), поэтому и становится итоговым победителем.

\clearpage

\section{Заключение}

В ходе выполнения данной лабораторной работы были освоены методы принятия решений на основе нечётких отношений предпочтения. Для варианта 4 были обработаны семь исходных отношений $R_1, \ldots, R_7$, заданных на множестве из шести альтернатив, и весовые коэффициенты $\lambda_1 = 0.1$, $\lambda_2 = 0.2$, $\lambda_3 = 0.1$, $\lambda_4 = 0.1$, $\lambda_5 = 0.3$, $\lambda_6 = 0.1$, $\lambda_7 = 0.1$.

Были построены два производных отношения. Пересечение $Q_1$ позволило выделить предпочтения, разделяемые всеми экспертами одновременно; после вычисления транспонированного отношения $Q_1^{*}$ и строгого отношения $Q_{1s}$ были найдены степени недоминируемости
\[
\mu_{\text{н.д.},Q_1} = (0.9,\; 0.8,\; 0.9,\; 0.7,\; 0.9,\; 1.0).
\]
Взвешенная свёртка $Q_2$ учла различную важность экспертов; для неё аналогичная процедура дала
\[
\mu_{\text{н.д.},Q_2} = (0.85,\; 0.67,\; 0.94,\; 0.86,\; 0.82,\; 0.87).
\]

Объединение двух подходов по правилу минимума дало итоговые степени недоминируемости
\[
\mu_{\text{н.д.}} = (0.85,\; 0.67,\; 0.9,\; 0.7,\; 0.82,\; 0.87).
\]
Наибольшее значение $0.9$ имеет альтернатива $x_3$, которая и признана оптимальной.

Полученный результат показывает, что использование только одного критерия может привести к ошибочному выбору: по пересечению лучшей выглядела альтернатива $x_6$, однако взвешенная свёртка показала её относительную слабость. Совместный анализ обоих отношений позволил выбрать альтернативу, устойчивую к изменению способа агрегирования экспертных мнений. Таким образом, цель лабораторной работы была полностью достигнута, а полученные навыки могут применяться при решении задач многокритериального выбора в условиях неполной и неточной информации.

\clearpage

\section{Список литературы}

\begin{enumerate}
    \item Заде Л. Понятие лингвистической переменной и его применение к принятию приближённых решений. --- М.: Мир, 1976. --- 165 с.
    \item Кофман А. Введение в теорию нечётких множеств. --- М.: Радио и связь, 1982. --- 432 с.
    \item Орловский С. А. Проблемы принятия решений при нечёткой исходной информации. --- М.: Наука, 1981. --- 208 с.
    \item Теория принятия решений в 2 т. Том 2: учебник и практикум для вузов / В. Г. Халин [и др.]; ответственный редактор В. Г. Халин. --- Москва: Издательство Юрайт, 2020. --- 431 с.
    \item Баранов В. В. Процессы принятия управляющих решений, мотивированных интересами: учебное пособие. --- Москва: ФИЗМАТЛИТ, 2005. --- 296 с.
    \item Лубенец Ю. В. Теория принятия решений: учебное пособие. --- Липецк: Изд-во ЛГТУ, 2011. --- 80 с.
\end{enumerate}
